\section*{ЗАКЛЮЧЕНИЕ}
\phantomsection
\addcontentsline{toc}{section}{Заключение}

В результате выполнения курсовой работы решена задача численного анализа влияния дискретного микрорельефа на трибологические характеристики радиального подшипника скольжения. Были выполнены следующие задачи:

\begin{enumerate}
    \item Проведён обзор теоретических основ гидродинамической смазки и современных исследований в области лазерного текстурирования поверхностей трения. Установлено, что ключевым механизмом положительного влияния микрорельефа является микрогидродинамический эффект, создаваемый локальными клиновыми зазорами на кромках углублений.

    \item Освоена методика численного решения уравнения Рейнольдса для подшипника конечной длины методом конечных разностей с итерационной процедурой Гаусса--Зейделя и верхней релаксацией. Расчёт выполнен на сетке $500 \times 500$ узлов с критерием сходимости $\delta < 10^{-5}$.

    \item Получены поля давления и интегральные характеристики (нагрузочная способность $F$, коэффициент трения $\mu$, расход смазки $Q$) для гладкого подшипника и подшипника с комбинированными двухуровневыми углублениями (тип~9) в диапазоне эксцентриситетов $\varepsilon = 0{,}05 - 0{,}80$. Геометрия подшипника: $R = 50$~мм, $c = 70$~мкм, $L = 80$~мм.

    \item Выполнен сравнительный анализ гладкого и текстурированного подшипников. При характерном рабочем эксцентриситете $\varepsilon = 0{,}6$ зафиксированы следующие изменения:
    \begin{itemize}
        \item нагрузочная способность увеличилась на $4{,}24$\%: с $F = 10685{,}6$~Н до $F = 11138{,}9$~Н;
        \item коэффициент трения снизился на $4{,}85$\%: с $\mu = 0{,}01507$ до $\mu = 0{,}01434$;
        \item расход смазки возрос на $0{,}48$\%: с $Q = 0{,}6739$~л/с до $Q = 0{,}6772$~л/с.
    \end{itemize}
\end{enumerate}

Полученные результаты подтверждают, что нанесение комбинированных двухуровневых углублений с параметрами $a = 2{,}48$~мм, $b = 2{,}28$~мм, $h_p = 10$~мкм ($h_p/c = 0{,}1$) приводит к улучшению рабочих характеристик подшипника. Одновременное повышение нагрузочной способности и снижение коэффициента трения при незначительном росте расхода смазки свидетельствует об эффективности данного типа микрорельефа.

Комбинированный двухуровневый профиль углублений, благодаря наличию плоской центральной зоны полной глубины, промежуточной ступени и плавного скругления у краёв, объединяет достоинства плоскодонного и криволинейного профилей. Плоская центральная зона обеспечивает стабильный резервуар смазки, промежуточная ступень формирует дополнительный клиновый эффект, а плавное скругление исключает резкие перепады давления на границах углубления. Это делает комбинированный профиль перспективным выбором для подшипников, работающих в стационарных режимах при умеренных и высоких эксцентриситетах.

На основании проведённого анализа можно рекомендовать применение комбинированного двухуровневого микрорельефа для подшипников скольжения нефтегазового оборудования, работающего в условиях жидкостного трения. Особенно перспективным является использование текстурирования в тяжелонагруженных узлах ($\varepsilon > 0{,}5$), где прирост нагрузочной способности и снижение трения выражены наиболее значительно.
